\section{Thesis Contributions}
\label{sec:}

Building upon the research objectives outlined above, this thesis contributes to the application of Artificial Intelligence in esports through the development of \textbf{a Human-in-the-Loop framework} with the possibility of performing \textbf{fully automated prediction of discursive functions} in \textit{League of Legends} esports shoutcasting in the future.

\medskip

Within this framework, an annotation tool with AI-assisted labeling will be implemented, enabling annotators to view model predictions when labeling difficult segments. This will ideally allow for a more efficient annotation process, while also providing insights on the process and possible improvements to the model.

\medskip

Therefore, we present the following list of contributions for this research project:
\begin{enumerate}
  \item \textbf{Annotation Framework}: a schema in order to annotate data with the discursive functions of shoutcasting for the development of predictive models and recommender systems.
  \item \textbf{Annotation Tool}: named \texttt{Momouchi-v2}, a web-based tool for the annotation of shoutcasting data, making use of the database of VODs available online, and making use of Youtube's automatic transcription of the shoutcaster's speech.
  \item \textbf{Annotated Dataset}: as a result of the two previous contributions, a computationally usable and structured dataset of shoutcasting segments labeled with their corresponding discursive functions.
  \item \textbf{Computational Modeling of Discursive Functions}: the development and evaluation of models for the automatic prediction of discursive functions in esports shoutcasting, with focus on traansformer-based architectures, and the use of context through adjacent segments and in-game events.
  \item \textbf{Human-in-the-Loop Evaluation}: a comprehensive evaluation of the developed framework and models, using the developed annotation tool. This will include user studies and a comparative analysis of the annotation process with and without assistance. 
\end{enumerate}

\medskip

If the proposed models achieve a high performance in our evaluation, it would be viable to allow it to perform fully automated labeling of the data. Consequently, this would facilitate the further analysis of discursive patterns in shoutcasting data, as well as having a larger dataset for the development of future tools, \textit{e.g.} a discourse-aware recommender system for shoutcasters. 

\medskip

In order to contextualize these contributions within the current state-of-the-art, the following chapter will review existing work more in depth. This includes prior approaches to the use of AI for esports analytics, discourse modeling, and shoutcasting. We aim to highlight the gaps in the current research that our work addresses, and further motivate the need for our proposed framework and models.

